% !PRINTED 365
\weberpage{365}
%<<CONT
gruppen auch durch $a'B$ darstellen, wo $a'$ in $A'$ enthalten ist. Sind dann $a_1'B, a_2'B$ \gesperrt{verschiedene} Nebengruppen zu $B$ in $A$, so sind $a_1'B', a_2'B'$ \gesperrt{verschiedene} Nebengruppen zu $B'$ in $A'$ und umgekehrt; denn $a_1'a_2'^{-1}$ ist dann und nur dann in $B'$ enthalten, wenn es zugleich in $B$ enthalten ist, und daraus folgt die Formel (9).

Diese S\"atze wenden wir nun auf die Ausdr\"ucke (2) und (3) f\"ur $h$ und $h'$ an. Nach (4) ist
\weq{10}{(\overline{O}, \overline{E}O') = (\overline{O}, \overline{E}O)(\overline{E}O, \overline{E}O'),}
also
\weq{11}{h' = h(\overline{E}O, \overline{E}O').}

Wir bezeichnen jetzt, wie schon oben, mit $E$ die Gruppe der numerischen Einheiten in $O$. Dann ist $\overline{E}E = \overline{E}$, da die numerischen Einheiten unter den funktionalen enthalten sind, und jede \gesperrt{Zahl}, die in $\overline{E}O'$ enthalten ist, ist in $EO'$ enthalten. Folglich ist $EO'$ der Durchschnitt von $O$ und $\overline{E}O'$. Wenden wir also (7) an, indem wir $O, EO', \overline{E}$ an Stelle von $A, B, C$ setzen, so ist $B$ der Durchschnitt von $A$ und $BC$ und wir erhalten:
\weq{12}{(\overline{E}O, \overline{E}O') = (\overline{E}O, \overline{E}EO') = (O, EO').}

Es ist ferner nach (4):
\weq{13}{(O, EO')(EO', O') = (O, O'),}
und wenn wir mit $E'$ die Gruppe der numerischen Einheit in $O'$ % [sic] Gedruckt ist Einheit im Singular., also den Durchschnitt von $E$ mit $O'$ bezeichnen, so da\ss{} $E'O' = O'$ ist:
\[
(EO', O') = (EO', E'O') = (E, E')
\]
(nach 2), denn $E'$ ist der Durchschnitt von $E'O'$ und $E$.

Also haben wir nach (13)
\[
(O, EO')(E, E') = (O, O')
\]
und nach (12)
\[
(\overline{E}O, \overline{E}O')(EE') = (O, O'), % [sic] Gedruckt ohne Komma zwischen E und E'.
\]
also schlie\ss lich nach (11):
\weq{14}{(E, E')h' = (O, O')h.}

15. \gesperrt{Satz. Die Formel (14) gilt unver\"andert, wenn $O'$ nicht gerade die Ordnungsgruppe, sondern irgend eine in $O$ enthaltene Zahlgruppe von endlichem Index $(O, O')$ bedeutet, wenn wir zwei Ideale $\mathfrak{a}, \mathfrak{a}'$ nach $O'$ \"aquivalent nennen, falls ihr Quotient $\mathfrak{a}'/\mathfrak{a} = \eta'$ eine Zahl in $O'$ ist.}
